\section*{Chapter \ref{ch:mx:order-types-and-their-uses} --- Order Types and Their Uses}

\begin{solution}{exo:mx:order-types-and-their-uses:1}
It is a \omterm{def:mx:order-types-and-their-uses:marketable}{marketable limit order}: it buys 300 at 100.01 and 300 at 100.02, leaving 200 offered at 100.02; nothing of it rests.
\end{solution}

\begin{solution}{exo:mx:order-types-and-their-uses:2}
The \omterm{def:mx:order-types-and-their-uses:tif}{fill-or-kill order} needs 400 and finds 300: it is accepted and cancelled whole, and nothing trades. The \omterm{def:mx:order-types-and-their-uses:tif}{immediate-or-cancel order} takes the 300
and cancels the remaining 100.
\end{solution}

\begin{solution}{exo:mx:order-types-and-their-uses:3}
Venues reward displayed liquidity: at one price, displayed orders rank ahead of non-displayed ones whatever their arrival times. A \omterm{def:mx:order-types-and-their-uses:hidden}{hidden order} buys
secrecy with priority.
\end{solution}

\begin{solution}{exo:mx:order-types-and-their-uses:4}
$10\,000/500=20$ slices appear, and the order goes to the back of the queue 19 times (every slice after the first).
\end{solution}

\begin{solution}{exo:mx:order-types-and-their-uses:5}
$\sigma\sqrt T=0.5\sqrt{30}=2.74$ ticks, $d/\sigma\sqrt T=0.365$: $C=2.74\,\varphi(0.365)-1\times(1-\Phi(0.365))=0.665$ ticks per share.
\end{solution}

\begin{solution}{exo:mx:order-types-and-their-uses:6}
A \omterm{def:mx:order-types-and-their-uses:stop}{stop-limit order} that reaches a book with no bids at or above its limit rests as an offer instead of selling lower, so it does not consume the next
bids or trigger the next stops: 25 cents instead of 37 with twenty stops. Its owner gives up certainty of getting out: in a real fall the order may
rest unexecuted while the price keeps going.
\end{solution}

\begin{solution}{exo:mx:order-types-and-their-uses:7}
\texttt{iceberg\_study} with \texttt{follow=True}: the strict detector catches 43\% of refills instead of 62\% (13 false alarms): a refill at the new ask
is not at the price of the slice that was consumed. The seller also sells only 23\,700 of the 30\,000 shares in the half hour, because refills that follow
a falling ask rest behind a new queue instead of meeting buyers at a price the market left.
\end{solution}

\begin{solution}{exo:mx:order-types-and-their-uses:8}
A \omterm{def:mx:order-types-and-their-uses:peg}{midpoint peg} saves half the spread only when it executes; it executes only when someone crosses to the mid, which is when that someone knows more
(adverse selection, chapter~9) or when a patient counterparty happens to be there. Orders that must complete take liquidity sometimes; a pegged-only
algorithm has execution risk and trades mostly against informed flow.
\end{solution}

\begin{solution}{pb:mx:order-types-and-their-uses:1}
\textbf{1.} Both rank behind displayed orders at their price for their hidden part; the iceberg shows a slice and its refreshes appear on the feed, the
\omterm{def:mx:order-types-and-their-uses:hidden}{hidden order} shows nothing.
\textbf{2.} A refresh is an increase of displayed size: it goes to the back of the queue with a new reference, as any new arrival.
\textbf{3.} From a partial fill that leaves an unwanted remainder or a partly hedged position; it costs the fills it forgoes when the book is short.
\textbf{4.} Each triggered stop becomes a market sell that consumes bids and prints lower prices, which trigger the next stops.
\textbf{5.} Five cents; 37 cents with twenty stops.
\textbf{6.} 19\,000 shares, for an initial sale of 3\,000.
\textbf{7.} 25 cents.
\textbf{8.} More depth between the triggers (liquidity providers stepping in), wider spacing of the stops, or \omterm{def:mx:order-types-and-their-uses:stop}{stop-limit orders}.
\textbf{9.} After an execution removes a displayed order entirely, flag an add on the same side and price within 50~ms whose size equals the removed
order's original size.
\textbf{10.} 99 refills, all 99 caught, with 20 false alarms.
\textbf{11.} Its refill is a new order that travels to the venue and meets the market as it is: when the price has moved, part of it executes on arrival
and the rest rests smaller, or at a different queue position; the venue's refresh appears in the same event, at the same size.
\textbf{12.} Detection falls to 18\% with the delay and 6\% with the size randomised too.
\textbf{13.} 47\% of refills (34 detections), with 378 false alarms: eleven false alarms for each true one.
\textbf{14.} \emph{Named result}: the strict detector catches 100\% of venue refreshes, 62\% of immediate algorithmic refills, 18\% with a 200~ms random
delay and 6\% with sizes randomised by 30\%; the loose detector recovers 47\% of the last at the price of 378 false alarms in half an hour.
\textbf{15.} The consumed slices' sizes over time (the same size repeating), the persistence of one \omterm{def:mx:the-limit-order-book:tick}{price level}, and execution imbalance at that level.
\textbf{16.} Time: a random delay leaves the ask without its liquidity, and slices sell later; in a moving market some are never filled.
\textbf{17.} Inferring hidden interest from public data is legal and common; trading on it is ordinary competition. What is not legal is obtaining it
from the venue or the broker (One Quant Book 9, chapter~29, for the enforcement record on manipulation; Book~16 for information barriers).
\textbf{18.} A \omterm{def:mx:order-types-and-their-uses:peg}{midpoint peg} in a dark venue, or an algorithm that varies slice size, price and timing.
\textbf{19.} Randomise the refreshed size within a band and delay the refresh by a random time; some venues offer exactly these options.
\textbf{20.} Its urgency, its size and its patience: every attribute is information for whoever reads the feed.
\end{solution}

\begin{solution}{iq:mx:order-types-and-their-uses:1}
When the trade must happen now or not at all: hedging, taking a price that will not last, sweeping several venues, pinging a dark pool. A resting \omterm{def:mx:the-limit-order-book:orders}{limit
order} reveals interest and risks being picked off. \iqlookfor{the link between order type and information leakage and adverse selection.}
\end{solution}

\begin{solution}{iq:mx:order-types-and-their-uses:2}
An order that becomes a \omterm{def:mx:the-limit-order-book:orders}{market order} when a trade prints at or through its stop price; it then executes against whatever the book holds, which in a fast
fall may be far lower (and other stops add to the selling). \iqlookfor{the conversion to a \omterm{def:mx:the-limit-order-book:orders}{market order}, the cascade, the stop-limit alternative.}
\end{solution}

\begin{solution}{iq:mx:order-types-and-their-uses:3}
A resting sell at $p$ is a call struck at $p$ written to the market: whoever learns first that value is above $p$ buys from it. A resting buy is a put.
The option expires when the owner cancels, so its value grows with the time the order is left unattended and with volatility. \iqlookfor{Copeland and
Galai; why cancellation speed is worth money.}
\end{solution}

\begin{solution}{iq:mx:order-types-and-their-uses:4}
Keep the displayed slice in the book and the reserve in the order record; when the slice is fully executed and reserve remains, move the next slice to
the back of the level with a new time stamp; on the feed, an execution of the old reference and an add under a new one (or a replace). Decide and
document the priority of the reserve against other orders. \iqlookfor{priority of the refreshed slice, and that the feed must not reveal the reserve.}
\end{solution}

\begin{solution}{iq:mx:order-types-and-their-uses:5}
Nasdaq's rule, for one: remove it from the book when the inside is crossed or missing and re-enter it at the new midpoint once there is a valid one,
cancelling it if none comes within a second. A midpoint defined by a locked or crossed quote is not a fair price. \iqlookfor{that pegged prices depend on
a reference that can be invalid, and the engine must say what happens then.}
\end{solution}

\begin{solution}{iq:mx:order-types-and-their-uses:6}
Look for refill patterns (the same size reappearing at the same price after executions), a level that absorbs far more than it displays, persistent
one-sided execution at one price, and hidden executions (non-displayed trade prints) clustering at a price; weigh each flag against its base rate of
coincidence. \iqlookfor{a concrete detector and an awareness of false alarms.}
\end{solution}
